% GENERATED FILE. Edit canonical lessons, then run npm run generate:books.
% canonical-source-hash: fnv1a64:191838eabc92dd8c
% canonical-lessons: PT-C319-santo, PT-C319-marcha, PT-C319-manjerico, PT-C319-sardinhada, PT-C319-bailarico

\chapter{Os santos populares}
\label{ch:pt-os-santos-populares}
\hlchaptermodality{319}

\begin{chapteropening}
\textbf{By the end of this chapter:} \emph{I can describe a June street festival: the saints, the parades, basil pots, sardines and dancing.}

Complete the last lesson of chapter 319: I can describe a June street festival: the saints, the parades, basil pots, sardines and dancing.
\end{chapteropening}

\section[o santo]{o santo — the saint}

\label{lesson:PT-C319-santo}



\begin{quote}
Before the new one: say the Portuguese for the toast, then the Portuguese for to toast.
\end{quote}

\subsection*{You'll want to know: o santo}
\textbf{o santo} — \textquotedblleft{}the saint\textquotedblright{}.

\begin{quote}
Cada aldeia tem o seu santo.
\end{quote}

\begin{quote}
\emph{Every village has its own saint.}
\end{quote}

Also an adjective, ‘holy’: \emph{a água santa}. Before a name it shortens to \emph{São}.

\begin{cousinweb}
\emph{santo} is inherited from Latin \emph{sānctus} \textquotedblleft{}holy, sacred\textquotedblright{} (EtymDB), from \emph{sancīre} \textquotedblleft{}to make sacred\textquotedblright{}.
\end{cousinweb}

\begin{grammarlens}[title={What you've built}]
One more word for telling, explaining and describing at B1.
\end{grammarlens}

\subsection*{Your turn}
\begin{itemize}
\raggedright
  \item \emph{Say it:} \emph{o santo}
  \item \emph{Say it:} \emph{o santo}, once more
  \item \emph{Say it:} a sentence of your own with \emph{o santo}
  \item \emph{Recall:} say the Portuguese for the toast, then the Portuguese for to toast, then say \emph{o santo} again
\end{itemize}

\begin{tcolorbox}[breakable,colback=teal!4,colframe=teal!35!black,title={Before you move on}]
What does o santo mean? (\textquotedblleft{}The saint\textquotedblright{}.) Say it once more.
\end{tcolorbox}
\section[a marcha]{a marcha — the march; a neighbourhood parade}

\label{lesson:PT-C319-marcha}



\begin{quote}
Before the new one: say the Portuguese for to toast, then the Portuguese for the saint.
\end{quote}

\subsection*{You'll want to know: a marcha}
\textbf{a marcha} — \textquotedblleft{}the march; a neighbourhood parade\textquotedblright{}.

\begin{quote}
A marcha do nosso bairro foi a mais bonita.
\end{quote}

\begin{quote}
\emph{Our neighbourhood’s parade was the most beautiful.}
\end{quote}

In Lisbon in June, neighbourhoods parade in \emph{marchas populares}, with songs and colourful clothes.

\begin{cousinweb}
\emph{marcha} came from French \emph{marche}, built on \emph{marcher} \textquotedblleft{}to walk\textquotedblright{}, which is usually traced to Frankish \emph{markōn} \textquotedblleft{}to mark, to tread\textquotedblright{} (EtymDB has no entry).
\end{cousinweb}

\begin{grammarlens}[title={What you've built}]
One more word for telling, explaining and describing at B1.
\end{grammarlens}

\subsection*{Your turn}
\begin{itemize}
\raggedright
  \item \emph{Say it:} \emph{a marcha}
  \item \emph{Say it:} \emph{a marcha}, once more
  \item \emph{Say it:} a sentence of your own with \emph{a marcha}
  \item \emph{Recall:} say the Portuguese for to toast, then the Portuguese for the saint, then say \emph{a marcha} again
\end{itemize}

\begin{tcolorbox}[breakable,colback=teal!4,colframe=teal!35!black,title={Before you move on}]
What does a marcha mean? (\textquotedblleft{}The march; a neighbourhood parade\textquotedblright{}.) Say it once more.
\end{tcolorbox}
\section[o manjerico]{o manjerico — the bush basil (in a pot)}

\label{lesson:PT-C319-manjerico}



\begin{quote}
Before the new one: say the Portuguese for the saint, then the Portuguese for the march.
\end{quote}

\subsection*{You'll want to know: o manjerico}
\textbf{o manjerico} — \textquotedblleft{}the bush basil (in a pot)\textquotedblright{}.

\begin{quote}
Ofereci um manjerico à minha namorada.
\end{quote}

\begin{quote}
\emph{I gave my girlfriend a pot of basil.}
\end{quote}

In June, a small pot of basil with a paper flower and a short poem is a typical gift.

\begin{cousinweb}
\emph{manjerico} is related to \emph{manjericão} \textquotedblleft{}basil\textquotedblright{}; the origin of both is not settled, and EtymDB has no entry.
\end{cousinweb}

\begin{grammarlens}[title={What you've built}]
One more word for telling, explaining and describing at B1.
\end{grammarlens}

\subsection*{Your turn}
\begin{itemize}
\raggedright
  \item \emph{Say it:} \emph{o manjerico}
  \item \emph{Say it:} \emph{o manjerico}, once more
  \item \emph{Say it:} a sentence of your own with \emph{o manjerico}
  \item \emph{Recall:} say the Portuguese for the saint, then the Portuguese for the march, then say \emph{o manjerico} again
\end{itemize}

\begin{tcolorbox}[breakable,colback=teal!4,colframe=teal!35!black,title={Before you move on}]
What does o manjerico mean? (\textquotedblleft{}The bush basil (in a pot)\textquotedblright{}.) Say it once more.
\end{tcolorbox}
\section[a sardinhada]{a sardinhada — the sardine feast}

\label{lesson:PT-C319-sardinhada}



\begin{quote}
Before the new one: say the Portuguese for the march, then the Portuguese for the bush basil.
\end{quote}

\subsection*{You'll want to know: a sardinhada}
\textbf{a sardinhada} — \textquotedblleft{}the sardine feast\textquotedblright{}.

\begin{quote}
Fizemos uma sardinhada no quintal.
\end{quote}

\begin{quote}
\emph{We had a sardine barbecue in the back yard.}
\end{quote}

Grilled sardines on bread, outdoors, typical of summer festivals.

\begin{cousinweb}
\emph{sardinhada} is built on \emph{sardinha}, from Latin \emph{sardīna} \textquotedblleft{}a sardine\textquotedblright{} (Lewis \& Short), with \emph{-ada} for a meal of something.
\end{cousinweb}

\begin{grammarlens}[title={What you've built}]
One more word for telling, explaining and describing at B1.
\end{grammarlens}

\subsection*{Your turn}
\begin{itemize}
\raggedright
  \item \emph{Say it:} \emph{a sardinhada}
  \item \emph{Say it:} \emph{a sardinhada}, once more
  \item \emph{Say it:} a sentence of your own with \emph{a sardinhada}
  \item \emph{Recall:} say the Portuguese for the march, then the Portuguese for the bush basil, then say \emph{a sardinhada} again
\end{itemize}

\begin{tcolorbox}[breakable,colback=teal!4,colframe=teal!35!black,title={Before you move on}]
What does a sardinhada mean? (\textquotedblleft{}The sardine feast\textquotedblright{}.) Say it once more.
\end{tcolorbox}
\section[o bailarico]{o bailarico — the street dance, a local dance party}

\label{lesson:PT-C319-bailarico}



\begin{quote}
Before the new one: say the Portuguese for the bush basil, then the Portuguese for the sardine feast.
\end{quote}

\subsection*{You'll want to know: o bailarico}
\textbf{o bailarico} — \textquotedblleft{}the street dance, a local dance party\textquotedblright{}.

\begin{quote}
Depois da sardinhada, houve bailarico na praça.
\end{quote}

\begin{quote}
\emph{After the sardine feast there was a dance in the square.}
\end{quote}

An informal word, from \emph{bailar} ‘to dance’ and \emph{baile}.

\begin{cousinweb}
\emph{bailarico} is built on \emph{bailar} \textquotedblleft{}to dance\textquotedblright{}, which goes back to Late Latin \emph{ballāre} \textquotedblleft{}to dance\textquotedblright{} (EtymDB; Lewis \& Short), with the ending \emph{-ico}.
\end{cousinweb}

\begin{grammarlens}[title={What you've built}]
One more word for telling, explaining and describing at B1.
\end{grammarlens}

\subsection*{Your turn}
\begin{itemize}
\raggedright
  \item \emph{Say it:} \emph{o bailarico}
  \item \emph{Say it:} \emph{o bailarico}, once more
  \item \emph{Say it:} a sentence of your own with \emph{o bailarico}
  \item \emph{Recall:} say the Portuguese for the bush basil, then the Portuguese for the sardine feast, then say \emph{o bailarico} again
\end{itemize}

\begin{tcolorbox}[breakable,colback=teal!4,colframe=teal!35!black,title={Before you move on}]
What does o bailarico mean? (\textquotedblleft{}The street dance, a local dance party\textquotedblright{}.) Say it once more.
\end{tcolorbox}
